\documentclass[11pt]{article}
\usepackage[a4paper,margin=1in]{geometry}
\usepackage{amsmath,amssymb,mathtools,bm}
\usepackage{enumitem,booktabs,array}
\usepackage{tcolorbox}
\usepackage{hyperref}
\usepackage[protrusion=true,expansion=false]{microtype}
\hypersetup{colorlinks=true,linkcolor=blue,urlcolor=blue,citecolor=blue}
\setlist[itemize]{topsep=2pt,itemsep=2pt}
\setlist[enumerate]{topsep=2pt,itemsep=2pt}
\newtcolorbox{keybox}{colback=blue!3!white,colframe=blue!40!black,boxrule=0.4pt,arc=1mm}
\newtcolorbox{alertbox}{colback=red!3!white,colframe=red!40!black,boxrule=0.4pt,arc=1mm}
\newtcolorbox{answerbox}{colback=green!3!white,colframe=green!40!black,boxrule=0.4pt,arc=1mm}
\newtcolorbox{drillbox}{colback=yellow!5!white,colframe=orange!60!black,boxrule=0.4pt,arc=1mm}
\newcommand{\EV}{\mathbb{E}}
\newcommand{\Var}{\mathrm{Var}}
\newcommand{\Cov}{\mathrm{Cov}}
\newcommand{\blank}[1]{\underline{\hspace{#1}}}

\title{W14-02: Duchin, Gao \& Xu (2024, JF) \\
\large ``Sustainability or Greenwashing? Evidence from the Asset Market for Industrial Pollution'' --- Active Recall Workbook}
\author{Banking \& Risk Management --- Exam Prep}
\date{\today}
\begin{document}\maketitle

\begin{keybox}
\textbf{Paper in one sentence.} When ESG-focused US firms \emph{divest} polluting industrial assets, the emissions don't fall --- they simply move to less ESG-sensitive buyers (often private/PE-owned or foreign) who operate the assets equally or more polluting --- so much of sell-side ``green'' activity is greenwashing, not actual emission reduction.

\textbf{Placement.} Canonical greenwashing paper; empirical complement to Hartzmark-Shue (2023) on ``green firms'' impact. A critical piece for the ESG-disclosure/regulation debate.
\end{keybox}

\section*{Round 1: Complete Walk-Through}

\subsection*{1.1 Design}
Asset-sale data on polluting plants (matched to EPA Toxics Release Inventory). Identify sellers (high-ESG-score public firms) and buyers (lower-ESG-score firms, private firms, or foreign firms). Compare post-sale plant emissions.

\subsection*{1.2 Empirical spec}
Plant-level panel DiD:
\[
\ln\text{Emissions}_{pt}=\alpha_p+\delta_t+\beta\cdot\text{PostSale}_{pt}+X_{pt}'\gamma+\varepsilon_{pt}.
\]
With plant and year FE, $\beta$ captures the change in emissions after ownership transfer from high- to low-ESG buyer.

\subsection*{1.3 Headline results}
\begin{itemize}
\item Post-sale, plant emissions stay the same or \emph{rise} modestly.
\item Buyers are disproportionately private-equity-owned, privately held, or foreign.
\item Sellers' corporate-level emission footprint falls mechanically; buyer reporting often does not.
\item Emissions effectively transferred across the market, not reduced.
\end{itemize}

\subsection*{1.4 Interpretation}
\begin{itemize}
\item Boundaries of the firm rather than production re-allocation drive measured ESG improvements.
\item ESG mandates may encourage divestment but not global-emission reduction.
\item Policy implication: emission-targeted regulation (carbon pricing, caps) more effective than firm-level ESG disclosure.
\end{itemize}

\subsection*{1.5 Caveats}
\begin{alertbox}
\begin{itemize}
\item Emissions post-sale may be under-reported by private buyers; true effect could be larger.
\item Sample of asset sales is limited; external validity to smaller transactions uncertain.
\item ``Greenwashing'' framing is interpretive; DGX provide the empirical result.
\end{itemize}
\end{alertbox}

\section*{Round 2: Level 1 Fill-in}
\begin{enumerate}
\item Asset type: polluting \blank{2cm} plants.
\item Emissions data source: EPA \blank{1cm}.
\item Post-sale emissions change: \blank{2cm}.
\item Typical buyer type: private, \blank{1cm}, or foreign.
\item Interpretation: \blank{1.5cm}.
\end{enumerate}

\section*{Round 3: Level 2 Prompts}
\begin{enumerate}
\item Write the plant-level DiD and discuss identification.
\item Why does sell-side ESG look green but aggregate emissions unchanged?
\item Discuss policy implications for ESG mandates vs.\ carbon taxes.
\item How could ESG rating systems incorporate these findings?
\end{enumerate}

\section*{Round 4: Minimal Scaffolding}
From a blank page: (i) asset-sale setting, (ii) DiD, (iii) post-sale emissions, (iv) interpretation, (v) caveats.

\section*{Round 5: Blank Page}
Essay: greenwashing in the market for industrial pollution.

\vspace{6pt}
\begin{drillbox}
\textbf{Speed Drill.} (1) Data source. (2) Typical buyers. (3) Emissions-change direction. (4) Policy implication. (5) Asset type.
\end{drillbox}

\section*{Quick Reference \& Answer Key}
\begin{answerbox}
\begin{itemize}
\item \textbf{Setting.} Asset sales of polluting plants.
\item \textbf{Effect.} Emissions unchanged or rise slightly.
\item \textbf{Buyers.} Private/PE/foreign, less ESG-sensitive.
\item \textbf{Lesson.} Firm-level ESG can be greenwashing.
\end{itemize}
\end{answerbox}

\begin{keybox}
\textbf{30-second exam pitch.} Duchin, Gao \& Xu (2024) use the asset market for industrial pollution to expose greenwashing. When ESG-focused public firms divest polluting plants, the plants are typically bought by private-equity, privately held, or foreign firms with less ESG pressure. Post-sale, plant emissions are unchanged or slightly higher. The seller's firm-level emissions footprint falls mechanically, but aggregate emissions are not reduced --- just transferred across the ownership boundary. The result is a powerful caution: firm-level ESG metrics are vulnerable to boundary manipulation, and ESG mandates that encourage divestment may yield misleading improvements without real emission reductions. Policy implication: emission-targeted regulation (cap-and-trade, carbon taxes) is more robust than firm-level ESG disclosure for driving real decarbonisation.
\end{keybox}

\end{document}
